% Generated by publication/build_sections.py; do not edit by hand.
\paragraph{Informal verdict.}
\textbf{A --- Complete solution}, with \textbf{high confidence} in the classification.

\paragraph{Lean coverage.}
Lean proves the complete classification of planar norms linearly equivalent to their duals, including the reflection and even-order rotational normal forms, automatic rotational symmetry, double duality, orthogonal classification, and power-bounded cosquare normal form.

\subsection{Audited informal answer}
Q805 --- complete classification
\begin{enumerate}
\item Faithful restatement
\end{enumerate}
Equip \(\mathbb R^2\) with its usual Euclidean inner product. For a norm \(N\) on \(\mathbb R^2\), define

\[
N^*(x)=\sup_{y\in\mathbb R^2\setminus\{0\}}
\frac{\langle x,y\rangle}{N(y)}.
\]

Characterize those norms \(N\) for which there is an invertible linear map \(u\in GL_2(\mathbb R)\) such that

\[
N^*=N\circ u.
\]

The PDF contains no subparts, footnotes, or additional qualifications. The absence of an absolute value in the displayed definition is harmless: replacing \(y\) by \(-y\) shows that the displayed supremum equals the usual dual norm with \(|\langle x,y\rangle|\). Luc Pommeret

\begin{enumerate}
\item Statement of the classification
\end{enumerate}
For \(\theta\in\mathbb R\), let \(R_\theta\) denote counterclockwise rotation through angle \(\theta\), and let

\[
F=\begin{pmatrix}1&0\\0&-1\end{pmatrix}
\]

be reflection in the first coordinate axis.
For a compact convex set \(K\subset\mathbb R^2\) containing \(0\) in its interior, write

\[
K^\circ=\{x\in\mathbb R^2:\langle x,y\rangle\le 1
\text{ for every }y\in K\}
\]

for its Euclidean polar.
Main theorem
Let \(N\) be a norm on \(\mathbb R^2\), and let

\[
B=\{x:N(x)\le 1\}
\]

be its unit ball. The following are equivalent.


There exists \(u\in GL_2(\mathbb R)\) such that \(N^*=N\circ u\).


There exists \(L\in GL_2(\mathbb R)\) such that, for \(C=LB\), one of the following holds:
Reflection type

\[
C=F C^\circ.
\tag{R}
\]

Rotational type

\[
C=R_{\pi/m}C^\circ
\tag{C\(_m\)}
\]

for some even integer \(m\ge 2\).


In case \((C_m)\), one automatically has

\[
R_{2\pi/m}C=C.
\tag{2.1}
\]

Conversely, every centrally symmetric convex body satisfying either \((R)\) or \((C_m)\) is the unit ball of a norm with the required property.
A norm may belong to both families; the theorem is an exhaustive classification, not necessarily a disjoint partition.
More explicitly, if \(C=LB\) and \(C=Q C^\circ\), where \(Q=F\) or \(Q=R_{\pi/m}\), then one may take

\[
\boxed{u=L^{-1}QL^{-T}.}
\tag{2.2}
\]

Thus every admissible norm can, after a linear change of coordinates, be made self-dual either by a reflection or by a half-step rotation relative to an even cyclic symmetry.

\begin{enumerate}
\item Translation into convex geometry
\end{enumerate}
We begin with elementary polar calculus.
Lemma 3.1 --- Polar identities
Let \(K\subset\mathbb R^2\) be compact, convex, and contain \(0\) in its interior. For \(A\in GL_2(\mathbb R)\) and \(r>0\),

\[
(AK)^\circ=A^{-T}K^\circ,
\qquad
(rK)^\circ=r^{-1}K^\circ,
\qquad
K^{\circ\circ}=K.
\tag{3.1}
\]

Proof
The first two identities follow directly from the definition. For example,

\[
x\in(AK)^\circ
\iff
\langle x,Ay\rangle\le1\quad(y\in K)
\]

which is equivalent to

\[
A^Tx\in K^\circ,
\]

hence \(x\in A^{-T}K^\circ\).
For the bipolar identity, \(K\subset K^{\circ\circ}\) is immediate. Let \(x\notin K\). Since \(K\) is compact, there is \(z\in K\) minimizing \(\|x-z\|_2\). Put \(p=x-z\neq0\). Convexity implies, for every \(y\in K\),

\[
\langle p,y-z\rangle\le0.
\]

Indeed, \(z+t(y-z)\in K\) for \(0\le t\le1\), and differentiating the squared distance from \(x\) at \(t=0\) gives this inequality. Consequently,

\[
\sup_{y\in K}\langle p,y\rangle
\le \langle p,z\rangle
<\langle p,x\rangle.
\]

Because \(0\) lies in the interior of \(K\),

\[
s:=\sup_{y\in K}\langle p,y\rangle>0.
\]

Then \(p/s\in K^\circ\), whereas

\[
\left\langle \frac p s,x\right\rangle>1.
\]

Thus \(x\notin K^{\circ\circ}\). Hence \(K^{\circ\circ}=K\). \(\square\)

Lemma 3.2 --- Unit balls
The unit ball of \(N^*\) is \(B^\circ\). The unit ball of \(N\circ u\) is \(u^{-1}B\). Consequently,

\[
N^*=N\circ u
\quad\Longleftrightarrow\quad
B^\circ=u^{-1}B
\quad\Longleftrightarrow\quad
B=uB^\circ.
\tag{3.2}
\]

Proof
By homogeneity,

\[
N^*(x)=\sup_{y\in B}\langle x,y\rangle.
\]

Thus \(N^*(x)\le1\) precisely when \(\langle x,y\rangle\le1\) for all \(y\in B\), namely when \(x\in B^\circ\).
The second assertion is immediate:

\[
N(ux)\le1 \iff ux\in B \iff x\in u^{-1}B.
\]

This proves (3.2). \(\square\)
So the problem is precisely:

Characterize the centrally symmetric convex bodies \(B\subset\mathbb R^2\) that are linearly equivalent to their polar.

The nontrivial part is to put the linear equivalence into canonical form.

\begin{enumerate}
\item A compactness constraint on the duality map
\end{enumerate}
Let

\[
\operatorname{Aut}(B)=\{A\in GL_2(\mathbb R):AB=B\}.
\]

Lemma 4.1
The group \(\operatorname{Aut}(B)\) is compact.
Proof
Choose \(0<r<R\) such that

\[
rD\subset B\subset RD,
\]

where \(D\) is the Euclidean unit disk. If \(A\in\operatorname{Aut}(B)\) and \(\|x\|_2\le1\), then \(rx\in B\), so

\[
rAx=A(rx)\in B\subset RD.
\]

Therefore

\[
\|A\|_{\mathrm{op}}\le \frac Rr.
\]

The same estimate applies to \(A^{-1}\), since \(A^{-1}B=B\). Hence the group is uniformly bounded and uniformly bounded away from singular matrices. It is closed in the space of \(2\times2\) matrices and therefore compact by Heine--Borel. \(\square\)
Now suppose that

\[
B=uB^\circ.
\tag{4.1}
\]

Taking polars gives

\[
B^\circ=u^{-T}B.
\]

Substitution into (4.1) yields

\[
B=uu^{-T}B.
\]

Thus the matrix

\[
S:=uu^{-T}
\tag{4.2}
\]

belongs to \(\operatorname{Aut}(B)\). In particular, all powers of \(S\) are bounded.
This seemingly modest observation almost completely determines the possible congruence type of \(u\).

\begin{enumerate}
\item Canonical form of an admissible duality
\end{enumerate}
Write

\[
H=\frac{u+u^T}{2},
\qquad
A=\frac{u-u^T}{2}.
\]

Thus \(H\) is symmetric and \(A\) is alternating. In dimension two,

\[
A=tJ,
\qquad
J=\begin{pmatrix}0&-1\\1&0\end{pmatrix}
\]

after a suitable congruence and scaling.
Recall that under a change \(u\mapsto TuT^T\),

\[
(TuT^T)(TuT^T)^{-T}
=
T(uu^{-T})T^{-1}.
\tag{5.1}
\]

Hence power-boundedness of \(uu^{-T}\) is preserved.
Also, replacing \(u\) by \(-u\) does not change \(uB^\circ\), since \(B^\circ=-B^\circ\).
Lemma 5.1 --- Bilinear normal form
Suppose \(u\in GL_2(\mathbb R)\) and \(uu^{-T}\) is power-bounded. Then there exist \(T\in GL_2(\mathbb R)\), \(c>0\), and \(Q\in O(2)\) such that

\[
Q=cTuT^T.
\tag{5.2}
\]

Moreover:


if \(\det u>0\), \(Q\) may be chosen to be a rotation;


if \(\det u<0\), \(Q\) may be chosen to be a reflection.


Proof
We distinguish the possible ranks and signatures of \(H\).
Case 1: \(H\) is definite
After replacing \(u\) by \(-u\) if necessary and applying a congruence, we may suppose

\[
H=I.
\]

Then

\[
u=I+tJ
=
\begin{pmatrix}
1&-t\\
t&1
\end{pmatrix}.
\]

Since

\[
(I+tJ)(I+tJ)^T=(1+t^2)I,
\]

the matrix

\[
Q=\frac{I+tJ}{\sqrt{1+t^2}}
\]

is an element of \(SO(2)\).

Case 2: \(H=0\)
Then \(u=tJ\) with \(t\neq0\). After multiplying by a positive scalar and possibly replacing \(u\) by \(-u\), this becomes \(J\), which is rotation through \(\pi/2\).

Case 3: \(H\) has rank one
After congruence and changing the sign of \(u\), we may write

\[
u=
\begin{pmatrix}
1&-t\\
t&0
\end{pmatrix},
\qquad t\neq0.
\]

A direct computation gives

\[
uu^{-T}
=
\begin{pmatrix}
-1&-2/t\\
0&-1
\end{pmatrix}.
\]

This is a nontrivial Jordan block with eigenvalue \(-1\), so its powers have off-diagonal entries growing linearly in magnitude. It is not power-bounded. Hence this case is impossible.

Case 4: \(H\) is indefinite
After congruence,

\[
u=
\begin{pmatrix}
1&-t\\
t&-1
\end{pmatrix},
\qquad t\neq\pm1.
\]

Here

\[
uu^{-T}
=
\frac1{1-t^2}
\begin{pmatrix}
1+t^2&2t\\
2t&1+t^2
\end{pmatrix}.
\tag{5.3}
\]

If \(t\neq0\), its trace is

\[
\frac{2(1+t^2)}{1-t^2}.
\]

Its absolute value is strictly larger than \(2\), so the two eigenvalues are real reciprocal numbers, one of absolute value greater than \(1\). Thus the powers are unbounded.
Consequently \(t=0\), and

\[
u=
\begin{pmatrix}
1&0\\
0&-1
\end{pmatrix}=F,
\]

an orthogonal reflection.
These four cases exhaust the possible symmetric parts \(H\), proving the lemma. \(\square\)

Corollary 5.2 --- Orthogonal reduction of the body
If \(B=uB^\circ\), then there are \(L\in GL_2(\mathbb R)\) and \(Q\in O(2)\) such that

\[
C:=LB
\qquad\text{satisfies}\qquad
C=QC^\circ.
\tag{5.4}
\]

Proof
Choose \(T,c,Q\) as in Lemma 5.1 and put

\[
L=\sqrt c\,T.
\]

Then, using \(B=uB^\circ\),

\[
QC^\circ
=
cTuT^T\left(c^{-1/2}T^{-T}B^\circ\right)
=
\sqrt c\,TuB^\circ
=
\sqrt c\,TB
=
C.
\]

\(\square\)
Thus it remains only to understand orthogonal self-polarities.

\begin{enumerate}
\item Classification of the orthogonal map
\end{enumerate}
Suppose that

\[
C=QC^\circ,
\qquad Q\in O(2).
\tag{6.1}
\]

Since \(Q\) is orthogonal,

\[
(QK)^\circ=QK^\circ.
\]

Taking polars in (6.1) gives

\[
C^\circ=QC.
\tag{6.2}
\]

Combining (6.1) and (6.2),

\[
Q^2C=C.
\tag{6.3}
\]

6.1. Orientation-reversing case
If \(\det Q=-1\), then \(Q\) is a Euclidean reflection. After an orthogonal change of coordinates it is \(F\). This gives precisely the reflection class

\[
C=FC^\circ.
\]

No further restriction follows, because \(F^2=I\).

6.2. Orientation-preserving case
Suppose that

\[
Q=R_\alpha.
\]

Equation (6.3) gives

\[
R_{2\alpha}C=C.
\tag{6.4}
\]

Also \(C=-C\), so \(R_\pi C=C\).
Let \(G\) be the subgroup of rotations generated by \(R_{2\alpha}\) and \(R_\pi\).
Infinite symmetry group
If \(G\) is infinite, it is dense in the circle group of rotations. Since \(C\) is closed and invariant under every element of \(G\), it is invariant under the closure of \(G\), hence under every Euclidean rotation.
Therefore \(C=rD\) for some \(r>0\). But \(C=QC^\circ\) then gives

\[
rD=\frac1rD,
\]

so \(r=1\). Thus \(C=D\).
The disk already belongs to \((C_2)\), since

\[
D=R_{\pi/2}D^\circ.
\]

Finite symmetry group
Suppose \(G\) is finite. Then for some even integer \(m\),

\[
G=\left\langle R_{2\pi/m}\right\rangle.
\tag{6.5}
\]

The integer \(m\) is even because \(R_\pi\in G\).
Put

\[
\delta=\frac{2\pi}{m}.
\]

Since \(Q^2=R_{2\alpha}\in G\), the class of \(\alpha\) modulo \(\delta\) satisfies

\[
2\alpha\equiv0\pmod\delta.
\]

Thus

\[
\alpha\equiv0
\quad\text{or}\quad
\alpha\equiv\frac\delta2
\pmod\delta.
\tag{6.6}
\]

If \(\alpha\equiv0\pmod\delta\), then \(Q\in G\), hence \(QC=C\). From \(C=QC^\circ\) we get \(C=C^\circ\).
But the only Euclidean self-polar convex body is the unit disk. Indeed, if \(C=C^\circ\), then every \(x\in C\) satisfies

\[
\|x\|_2^2=\langle x,x\rangle\le1,
\]

so \(C\subset D\). Taking polars gives

\[
D=D^\circ\subset C^\circ=C,
\]

hence \(C=D\).
In the other case, choose \(H\in G\) such that

\[
HQ=R_{\delta/2}=R_{\pi/m}.
\]

Since \(HC=C\), equation \(C=QC^\circ\) becomes

\[
C=R_{\pi/m}C^\circ.
\]

This is exactly the rotational normal form \((C_m)\). Moreover,

\[
R_{2\pi/m}C=C
\]

follows from (6.3).
This proves the main theorem.

\begin{enumerate}
\item A completely explicit radial-function characterization
\end{enumerate}
The preceding theorem reduces the problem to two canonical functional equations.
For a centrally symmetric convex body \(C\), let

\[
e_\theta=(\cos\theta,\sin\theta),
\]

and define its radial and support functions by

\[
\rho(\theta)=\max\{r\ge0:re_\theta\in C\},
\]


\[
h(\theta)=\max_{x\in C}\langle e_\theta,x\rangle.
\]

The radial function of the polar is

\[
\rho_{C^\circ}(\theta)=\frac1{h(\theta)}.
\tag{7.1}
\]

Indeed, \(te_\theta\in C^\circ\) exactly when

\[
t\,h(\theta)\le1.
\]

If \(C=QC^\circ\), then \(C^\circ=Q^TC\), and therefore

\[
\frac1{h(\theta)}
=
\rho_{C^\circ}(\theta)
=
\rho_C(Qe_\theta).
\tag{7.2}
\]

Hence:
Reflection type

\[
\boxed{h(\theta)\rho(-\theta)=1
\quad\text{for every }\theta.}
\tag{7.3}
\]

Rotational type
For some even \(m\ge2\),

\[
\boxed{
h(\theta)\rho\!\left(\theta+\frac{\pi}{m}\right)=1
\quad\text{for every }\theta.
}
\tag{7.4}
\]

Equation (7.4) automatically implies

\[
\rho\!\left(\theta+\frac{2\pi}{m}\right)=\rho(\theta).
\]

Conversely, either (7.3) or (7.4) implies the corresponding equality of bodies, because two star-shaped bodies containing the origin are equal when their radial functions agree.
This gives a literal one-variable parametrization. Namely, begin with a positive continuous \(\pi\)-periodic function \(\rho\) such that

\[
C_\rho=\{re_\theta:0\le r\le\rho(\theta)\}
\]

is convex. Compute

\[
h_\rho(\theta)
=
\max_{\varphi\in\mathbb R}
\rho(\varphi)\cos(\theta-\varphi).
\]

Then \(C_\rho\) gives a solution exactly when, after a linear change of coordinates, either

\[
h_\rho(\theta)\rho(-\theta)=1,
\]

or, for an even \(m\),

\[
h_\rho(\theta)
\rho\!\left(\theta+\frac{\pi}{m}\right)=1.
\]


\begin{enumerate}
\item Examples and boundary cases
\end{enumerate}
8.1. Ellipsoidal norms
Let \(A\) be symmetric positive definite and

\[
N_A(x)=\sqrt{x^TAx}.
\]

Then

\[
N_A^*(x)=\sqrt{x^TA^{-1}x}.
\]

Taking

\[
u=A^{-1}
\]

gives

\[
N_A(ux)
=
\sqrt{x^TA^{-1}x}
=
N_A^*(x).
\]

After a linear change of coordinates, the unit ball becomes the Euclidean disk.

8.2. Regular even polygons
Let \(m\ge4\) be even and let \(P_m(r)\) be the regular \(m\)-gon with vertices

\[
r e_{2\pi k/m},\qquad 0\le k<m.
\]

Its polar has vertices in the directions

\[
\frac{(2k+1)\pi}{m}
\]

and circumradius

\[
\frac1{r\cos(\pi/m)}.
\]

Choose

\[
r=\frac1{\sqrt{\cos(\pi/m)}}.
\]

Then the polar has the same circumradius, and

\[
P_m(r)^\circ=R_{\pi/m}P_m(r).
\]

Since \(P_m(r)\) is invariant under \(R_{2\pi/m}\), it also satisfies

\[
P_m(r)=R_{\pi/m}P_m(r)^\circ.
\]

Thus every even \(m\ge4\) occurs.
The case \(m=2\) is the quarter-turn relation

\[
C=R_{\pi/2}C^\circ.
\]

These are traditionally the planar bodies self-polar with respect to an alternating area form; the Euclidean disk and a suitably scaled regular hexagon are examples.

8.3. The square and the \(\ell_\infty\) norm
For

\[
N(x,y)=\max(|x|,|y|),
\]

the dual norm is

\[
N^*(x,y)=|x|+|y|.
\]

With

\[
u=
\begin{pmatrix}
1&-1\\
1&1
\end{pmatrix}
=
\sqrt2\,R_{\pi/4},
\]

we have

\[
N(u(x,y))
=
\max(|x-y|,|x+y|)
=
|x|+|y|.
\]

Thus \(N^*=N\circ u\). This is the \(m=4\) rotational case before normalization of the scalar factor.
For \(N=\|\cdot\|_1\), one may take

\[
u=\frac12
\begin{pmatrix}
1&1\\
1&-1
\end{pmatrix}.
\]


8.4. A genuinely reflection-type example
Let

\[
v_0=\left(1,\frac15\right),\qquad
v_1=(1,0),\qquad
v_2=\left(\frac34,-\frac54\right),
\]

and put

\[
v_{i+3}=-v_i\qquad(0\le i\le2).
\]

Let

\[
K=\operatorname{conv}\{v_0,\ldots,v_5\}.
\]

The polar vertices corresponding to the first three consecutive edges are

\[
w_0=(1,0),\qquad
w_1=\left(1,-\frac15\right),\qquad
w_2=\left(-\frac34,-\frac54\right),
\]

and the remaining three are their negatives. These are exactly the images under \(F\) of the six vertices of \(K\). Hence

\[
K^\circ=FK
\]

and therefore

\[
K=FK^\circ.
\]

This example does not possess an orientation-preserving duality. Indeed, the absolute determinants of consecutive vertex pairs are

\[
\frac15,\quad \frac54,\quad \frac75,\quad
\frac15,\quad \frac54,\quad \frac75.
\tag{8.1}
\]

A linear automorphism of \(K\) must permute the vertices cyclically or anticyclically and has determinant of absolute value \(1\), since it preserves the area of \(K\). The cyclic word in (8.1), whose three entries are distinct, has only the two rotational symmetries corresponding to the identity and central inversion and has no reflection symmetry. Consequently,

\[
\operatorname{Aut}(K)=\{\pm I\}.
\]

If \(vK^\circ=K\), then, since \(FK^\circ=K\),

\[
vF^{-1}K=K.
\]

Thus \(vF\in\operatorname{Aut}(K)\), so \(v=\pm F\). Every duality has negative determinant. This shows that the reflection branch cannot be omitted.

8.5. The classical \(\ell_p^2\) norms
Let \(1<p<\infty\), and let \(q=p/(p-1)\). The dual of \(\ell_p^2\) is \(\ell_q^2\).
The norm \(\ell_p^2\) has the required property exactly when \(p=2\).
Indeed, if \(p\neq2\), one of \(p,q\) is greater than \(2\) and the other lies between \(1\) and \(2\). Suppose \(p>2>q\).
Near \((1,0)\), the boundary of the \(p\)-ball is

\[
x=(1-|y|^p)^{1/p}.
\]

For \(p>2\), this is twice continuously differentiable at \(y=0\). The corresponding boundary of the \(q\)-ball is

\[
x=(1-|y|^q)^{1/q}
=
1-\frac1q|y|^q+o(|y|^q),
\]

and

\[
\frac{x(y)-x(0)-x'(0)y}{y^2}
\sim -\frac1q |y|^{q-2}\longrightarrow-\infty.
\]

Thus the \(q\)-sphere is not \(C^2\) at the coordinate axes, whereas the \(p\)-sphere is \(C^2\) everywhere. An invertible linear map preserves the property of being a \(C^2\) embedded curve, so the two balls cannot be linearly equivalent.
Hence, among \(1<p<\infty\), only \(p=2\) works. The endpoint norms \(p=1,\infty\) do work, as shown above.
This also demonstrates that smoothness and strict convexity alone are far from sufficient.

\begin{enumerate}
\item Polygonal decision algorithm
\end{enumerate}
For a polygonal norm the property is exactly decidable using finite linear algebra.
Assume that the nonredundant vertices of \(B\) are

\[
v_0,\ldots,v_{2n-1}
\]

in cyclic order.


For each \(i\), compute the polar vertex \(w_i\) by solving

\[
\langle w_i,v_i\rangle=1,
\qquad
\langle w_i,v_{i+1}\rangle=1.
\]



Any linear map carrying \(B^\circ\) onto \(B\) must carry the cyclically ordered polar vertices either cyclically or anticyclically onto the vertices of \(B\).


For each sign \(\varepsilon\in\{1,-1\}\) and each shift \(k\), define the unique linear map \(u\) satisfying

\[
uw_0=v_k,\qquad
uw_1=v_{k+\varepsilon}.
\]



Verify whether

\[
uw_i=v_{k+\varepsilon i}
\]

for every \(i\), with indices modulo \(2n\).


There are \(4n\) candidates, and each is checked using \(O(n)\) exact arithmetic operations. Thus the naive complexity is \(O(n^2)\). If the vertices are rational, the entire test can be carried out over \(\mathbb Q\).
No computation is used in the proof of the classification theorem.

\begin{enumerate}
\item Topological and logical assumptions
\end{enumerate}
The proof uses only finite-dimensional Euclidean topology.
Separation
The identity \(K^{\circ\circ}=K\) uses finite-dimensional strict separation of a point from a compact convex set. It was proved directly using a nearest point, so no Hahn--Banach theorem is required.
Compactness
Compactness is used in three places:


the unit ball of a norm is compact by Heine--Borel;


support functions attain their maxima;


\(\operatorname{Aut}(B)\) is compact, which forces \(uu^{-T}\) to be power-bounded.


No weak topology or Banach--Alaoglu theorem occurs.
Countability and sequential arguments
No countability assumption is imposed on any family. The only sequence considered is the countable sequence of powers of a rotation. When this subgroup is dense, invariance passes to its closure by taking convergent sequences. This is legitimate because \(SO(2)\) and \(\mathbb R^2\) are metrizable; no nets or general nonsequential compactness arguments are involved.
Choice
No form of the axiom of choice beyond finite choices is needed. There is no Zorn lemma, no selection of representatives from an infinite family, and no nonconstructive compactness principle.
Smoothness and strict convexity
Neither is assumed. The regular polygon examples show that nonsmooth and non-strictly-convex norms occur naturally.
Invertibility of \(u\)
The explicit hypothesis \(u\in GL_2(\mathbb R)\) is actually redundant. If a linear map \(u\) satisfied \(N^*=N\circ u\), then \(ux=0\) would imply \(N^*(x)=0\), hence \(x=0\). Thus \(u\) is injective and therefore invertible in dimension two.
Dependence on dimension
The reduction of the symmetric part of \(u\) into the four cases above is specifically two-dimensional. In higher dimensions the congruence classification of bilinear forms and the possible compact ``cosquares'' \(uu^{-T}\) are substantially richer.

\begin{enumerate}
\item Final form of the answer
\end{enumerate}
A norm \(N\) on \(\mathbb R^2\) is linearly isometric to its dual precisely when, after a nonsingular linear change of coordinates, its unit ball \(C\) is one of the following:

\[
\boxed{C=FC^\circ}
\]

for a reflection \(F\), or

\[
\boxed{C=R_{\pi/m}C^\circ}
\]

for some even integer \(m\ge2\); in the latter case \(C\) necessarily has \(m\)-fold rotational symmetry.
Equivalently, in radial/support coordinates, the complete alternatives are

\[
\boxed{h(\theta)\rho(-\theta)=1}
\]

or

\[
\boxed{
h(\theta)\rho\!\left(\theta+\frac{\pi}{m}\right)=1
\quad(m\ge2\text{ even}).
}
\]

These conditions are both necessary and sufficient.
